\section{Moment reconstruction with an estimated geometric factor}
\label{sec:moment-factor-details}
The finite law theorem uses three different approximations: a geometric
factor is reconstructed, occupation moments are sampled, and a positive
probability is fitted to those moments.  We give quantitative interfaces
between these operations.  In particular, the fitted geometric moments
can be the exact moments of a rational polygon.  They then belong to an
actual probability measure at every degree used in the reconstruction.

Fix the complete component used in
Theorem~\ref{thm:two-field-finite-law}.  The constants in the lemmas below
depend only on a bound for its aperture, a lower bound $g_*>0$ for the
gap between expanded components, a lower bound $v_*>0$ for the obstacle
area, an upper bound $p_*$ for the relevant convex perimeters, and the
fixed prefix length $M$.  These bounds are supplied by the quantitative
geometric priors.  All assertions about estimated geometry are made on
the geometric stage's common accuracy event.  Numerical integration and
rational arithmetic are deterministic operations conditional on that
event.

Enlarge the fixed rational scale $R$ of
Section~\ref{sec:two-field-finite-law}, if necessary, so that the selected
body, the centered footprint, and their small approximation collars lie
in $(-R/4,R/4)^2$.  The cutoff support lies in $(-R/2,R/2)^2$.
This is a choice of coordinates within the fixed aperture.  We shall use
candidate normalized displacements in $[-1/2,1/2]^2$; consequently all
coordinate differences appearing in the factorization have absolute
value less than one.  For a finite signed measure, $\|\cdot\|_{\rm TV}$
denotes its full variation norm.

\subsection{A complete component and a rational geometric factor}
\begin{lemma}[A protected cutoff with rational values]
\label{lem:moment-protected-cutoff}
Let $P=C_0+(-A_0)$ be a complete expanded component, and let
$\widehat P$ be a rational convex polygon with
$d_H(\widehat P,P)\le\xi$.  Choose a fixed positive rational $r$ with
$r<g_*/(16\sqrt2)$ and suppose $\xi<r$.  There is an explicitly
specified cutoff $H$, with $0\le H\le1$, which is one on $P$, zero
on every other expanded component, and has Lipschitz constant at most
$r^{-1}$.  Its support is contained in
$\widehat P+[-2r,2r]^2$, and its values at rational points are rational.
For every compact launch probability with support $A_0$, including a
singular probability, one has the pointwise identity
\begin{equation}\label{eq:moment-cutoff-identity}
 H(x)v(x)=v_C(x)=\int\one_{C_0}(x+z)\,\dd\mu_0(z).
\end{equation}
\end{lemma}
\begin{proof}
Let $d_\infty(x,\widehat P)$ be distance for the maximum norm, and put
\[
 \chi(u)=
 \begin{cases}
 1,&0\le u\le1,\\
 2-u,&1<u<2,\\
 0,&u\ge2,
 \end{cases}
 \qquad
 H(x)=\chi\bigl(d_\infty(x,\widehat P)/r\bigr).
\]
Distance to a closed set is one-Lipschitz.  The bound on $\xi$ gives
$H=1$ on $P$.  If $P'$ is another expanded component, then
\[
 \inf_{x\in P'}d_\infty(x,\widehat P)
       \ge g_*/\sqrt2-\xi>2r,
\]
so $H$ is zero there, including the boundary.  The support and
Lipschitz assertions follow directly from the definition.  Distance
from a rational point to a rational polygon in the maximum norm is
the optimum of a finite rational linear programme; its optimum is
rational.  Thus $H$ has the asserted rational values.

For any component $C'$, its occupation contribution vanishes outside
$C'+(-A_0)$: if $x$ is outside this Minkowski sum, $x+z$ is outside
$C'$ for every $z\in A_0$.  The local finiteness and compact support
make the sum of these contributions locally finite.  Multiplication
by $H$ therefore retains precisely the contribution of $C_0$.
This argument is pointwise and does not discard boundary atoms.
\end{proof}

\begin{lemma}[Rational polygon factors and normalization]
\label{lem:moment-rational-factor}
Suppose the recovered convex body $\widehat C$ has a finite
computable support representation and
$d_H(\widehat C,C_0)\le\xi$.  For every positive rational $\zeta$
one can construct a rational convex polygon $D\supset\widehat C$
with $d_H(D,\widehat C)\le\zeta$.  Put $\tau=\xi+\zeta$ and assume
it is sufficiently small.  Then $|D|\ge v_*/2$ and
\begin{equation}\label{eq:moment-factor-TV}
 \left\|\frac{\one_{C_0}}{|C_0|}\dd x
              -\frac{\one_D}{|D|}\dd x\right\|_{\rm TV}
       \le C\tau.
\end{equation}
All normalized moments
\begin{equation}\label{eq:moment-rational-factor-moments}
 c_\alpha=\frac1{|D|}\int_D(x/R)^\alpha\dd x
\end{equation}
are rational and can be computed exactly by finite arithmetic;
$c_0=1$.  Their discrepancies from the corresponding moments of
the uniform law on $C_0$ are at most $C\tau$, independently of
degree.

Suppose rational estimates $\widehat I_\alpha$ of the occupation
moments have errors at most $\eta$.  Define
\begin{equation}\label{eq:moment-rational-data}
 \widetilde y_0=1,\qquad
 \widetilde y_\alpha=\widehat I_\alpha/|D|
                              \quad(\alpha\ne0).
\end{equation}
If $X$ is uniform on $C_0$ and independent of $Z\sim\mu_0$, then
\begin{equation}\label{eq:moment-normalization-error}
 \left|\widetilde y_\alpha-
                   \E\bigl((X-Z)/R\bigr)^\alpha\right|
       \le C(\eta+\tau).
\end{equation}
\end{lemma}
\begin{proof}
Here is a finite outer-polygon construction.  Take a sufficiently
fine net of rational normal vectors, including the four coordinate
directions.  Compute rational upper bounds for the support of
$\widehat C$ at these vectors, with a prescribed rational error,
and intersect the corresponding half-planes.  The coordinate
directions bound this intersection in a fixed square.  If $q$ is
within $\omega$ of a unit vector $u$, then, for every point $x$ in
that square,
\[
 u\cdot x\le q\cdot x+C\omega,
 \qquad
 |h_{\widehat C}(q)-h_{\widehat C}(u)|\le C\omega.
\]
Consequently, normal mesh and support-value errors at most
$c\zeta$ imply $h_D(u)\le h_{\widehat C}(u)+\zeta$ for every $u$.
The reverse inequality follows from containment.  The vertices of
this bounded intersection are rational.  The finite geometric
representation supplies support values with certified error, so all
operations in this construction terminate.

For convex bodies at Hausdorff distance at most $\tau$, the two
containments in their $\tau$-parallel bodies and the planar
parallel-body formula give
\[
 |C_0\mathbin\triangle D|
       \le (\operatorname{Per}(C_0)+\operatorname{Per}(D))\tau
                    +2\pi\tau^2
       \le 2p_*\tau+2\pi\tau^2.
\]
The relevant perimeters have a uniform bound because the bodies
are convex and lie in a fixed ball.  This also proves
$|D|\ge v_*/2$ for sufficiently small $\tau$.  Directly comparing
the two normalized indicator densities yields
\[
 \left\|\frac{\one_{C_0}}{|C_0|}
              -\frac{\one_D}{|D|}\right\|_1
 \le \frac{2|C_0\mathbin\triangle D|}
                       {\min\{|C_0|,|D|\}},
\]
which proves \eqref{eq:moment-factor-TV}.  A normalized coordinate
monomial is bounded by one on both bodies.  Integration against the
signed difference of their uniform probabilities therefore gives
the degree-independent moment estimate.

Triangulate $D$ using its rational vertices.  An affine map from
the reference triangle converts each monomial integral into a finite
sum of rational multiples of
\[
 \int_{\substack{s,t\ge0\\s+t\le1}}s^p t^q\dd s\dd t
                       =\frac{p!q!}{(p+q+2)!}.
\]
Its Jacobian, coefficients, and the scale $R$ are rational.  This
computes both the area and every integral in
\eqref{eq:moment-rational-factor-moments} exactly.  In particular,
the resulting array is the moment array of the actual uniform
probability on $D$.

Finally, Fubini gives
$I_\alpha=|C_0|\E((X-Z)/R)^\alpha$.  The scale is in the
monomial, so this identity has the physical area $|C_0|$ as its
normalizing factor.  Since $|I_\alpha|\le|C_0|$,
\[
 \left|\frac{\widehat I_\alpha}{|D|}
                    -\frac{I_\alpha}{|C_0|}\right|
 \le \frac{\eta+\bigl||D|-|C_0|\bigr|}{|D|}
 \le C(\eta+\tau).
\]
For degree zero the assigned value is exact.  This proves
\eqref{eq:moment-normalization-error}.
\end{proof}

\subsection{Shared observations and the conditioning bound}
\begin{lemma}[Uniform quadrature for all the occupation moments]
\label{lem:moment-shared-quadrature}
Conditional on the cutoff in Lemma~\ref{lem:moment-protected-cutoff},
let $n\ge1$, $0<\eta<1/4$, and $0<\delta<1/4$.  Fresh observations
with the two fixed commands give rational estimates of all
$I_\alpha$, $|\alpha|\le n$, with simultaneous error at most
$\eta$ and conditional probability at least $1-\delta$, using at
most
\begin{equation}\label{eq:moment-shared-quadrature-cost}
 C(n+1)^2\eta^{-4}
                   \log\frac{C(n+1)}{\eta\delta}
\end{equation}
attempted bits.  The nominal mesh is dyadic and comparable to
$\eta/(n+1)$.  The same empirical occupation values serve every
multi-index.  This conclusion is uniform over all compact launch
probabilities with the stated support bound.
\end{lemma}
\begin{proof}
Tile a fixed square containing the support of $H$, with a fixed
collar, by squares of side $\ell$.  Use one fixed grid corner
$x_Q$ of each square $Q$.  There are at most $C\ell^{-2}$ such
points.  For fixed $z$, put
\[
 w_\alpha(x)=H(x)(x/R)^\alpha,
 \qquad
 f_{\alpha,z}(x)=w_\alpha(x)\one_{C_0}(x+z).
\]
On the square under consideration, the weight is bounded by one
and has Lipschitz constant at most $C(n+1)$.  The union of the
grid squares that meet $\partial(C_0-z)$ has area at most
$C\ell$, uniformly in $z$.  Indeed, it lies in the
$\sqrt2\ell$-neighborhood of this convex boundary, whose area is
bounded by a constant times its perimeter times $\ell$, plus
$C\ell^2$.  On each remaining square the indicator is constant,
and the weight varies by at most $C(n+1)\ell$.  It follows that
\begin{equation}\label{eq:moment-fixed-z-quadrature}
 \left|\ell^2\sum_Q f_{\alpha,z}(x_Q)
                   -\int f_{\alpha,z}(x)\dd x\right|
                       \le C(n+1)\ell.
\end{equation}
Values on a closed obstacle boundary occur only in the boundary
squares already counted.  Thus the estimate holds for every $z$,
and integration against any $\mu_0$ preserves it.  This establishes
the required quadrature bound without regularity of the occupation
field.

At every $x_Q$, use the $2M$ means in the finite-prefix formula.
Estimate each of these means to accuracy $\sigma/(2M)$ with fresh
Bernoulli observations.  The Lipschitz bound for the prefix maximum
then gives $|\widehat v(x_Q)-v(x_Q)|\le\sigma$; truncation to
$[0,1]$ preserves this bound.  Hoeffding's inequality and a union
bound over the at most $2MC\ell^{-2}$ means require at most
\[
 C\ell^{-2}\sigma^{-2}\log\frac{C}{\ell\delta}
\]
attempted bits.  Constants absorb the fixed powers of $M$.
On the resulting event, replacing $v(x_Q)$ by $\widehat v(x_Q)$
in any weighted sum changes it by at most $C\sigma$.
Choose $\ell$ dyadic and comparable to $c\eta/(n+1)$ and set
$\sigma=c\eta$, with $c$ sufficiently small.  Together with
\eqref{eq:moment-fixed-z-quadrature}, this proves the simultaneous
accuracy and \eqref{eq:moment-shared-quadrature-cost}.  All sums
are rational.  The confidence bound is conditional on earlier data
because this stage uses fresh preparations.  The event controlling
the occupation values is common to every monomial, so the number
of moments does not multiply the observation bound.
\end{proof}

\begin{lemma}[An explicit amplification factor]
\label{lem:moment-explicit-conditioning}
Under the hypotheses of Lemma~\ref{lem:two-field-moment-separation},
with moment discrepancy bounded by $b$ through degree $4m$, one has
\begin{equation}\label{eq:moment-explicit-amplification}
 W_1(\mathcal L(Z),\mathcal L(\widetilde Z))
       \le \frac{C}{m}+\Lambda_m b,
 \qquad
 \Lambda_m=C_*m^2 24^{4m}(4m)!,
\end{equation}
where $C_*$ is numerical.  In particular
$\Lambda_m\le(Cm)^{4m+2}$ after increasing the numerical $C$.
\end{lemma}
\begin{proof}
The triangular identity
\eqref{eq:two-field-triangular-moments} gives, for the largest
degree-$k$ discrepancy $D_k$ in the unknown factor,
\[
 D_0=0,\qquad
 D_k\le b(1+2^k)+\sum_{j=0}^{k-1}\binom{k}{j}D_j.
\]
The coefficient of the unknown degree-$k$ moment has modulus one;
all moments of the two individual factors are bounded by one.
Induction gives $D_k\le2b8^k k!$.  To check the induction
quantitatively, the lower-degree sum divided by $2b8^k k!$ is
at most
\[
 \sum_{j=1}^k\frac{8^{-j}}{j!}
            \le\sum_{j=1}^\infty8^{-j}=1/7,
\]
whereas the forcing term divided by the same bound is at most
$1/4$ for $k\ge1$.  The case $b=0$ follows directly from the
triangular identity.

The tensor polynomial construction in the proof of
Lemma~\ref{lem:two-field-moment-separation} approximates any
one-Lipschitz function, after subtraction of its value at zero,
with uniform error at most $C/m$.  Its total degree is at most
$4m$, and the sum of its absolute monomial coefficients is at
most $Cm^2 3^{4m}$.  Applying the preceding moment bound gives
an integration error at most
$Cb m^2 3^{4m}8^{4m}(4m)!$.  Transportation duality proves
\eqref{eq:moment-explicit-amplification}.  Finally
$(4m)!\le(4m)^{4m}$ gives the stated elementary upper bound for
$\Lambda_m$.
\end{proof}

\subsection{Positive reconstruction and its finite output}
\begin{lemma}[A positive rational programme with a feasible comparison law]
\label{lem:moment-positive-programme}
Let $D$ and $c_\alpha$ be as in
Lemma~\ref{lem:moment-rational-factor}, and let $K_A$ be a rational
polygon with $d_H(K_A,A_0)\le\rho$, where $\rho$ is a known
positive rational bound on the geometric accuracy event.  For a
sufficiently small positive dyadic $h$, retain the grid
\begin{equation}\label{eq:moment-padded-grid}
 G=\left\{g\in h\Z^2\cap[-1/2,1/2]^2:
   d_\infty(g,R^{-1}K_A)\le\rho/R+2h\right\}.
\end{equation}
It is nonempty and can be computed by rational arithmetic.  Every
point of $R^{-1}A_0$ has a retained grid point within $\sqrt2h$.

Put $n=4m$, and let the rational numbers $\widetilde y_\alpha$
have errors at most $e_y$ as moments of $Y=(X-Z)/R$.  For a
probability vector $p=(p_g)_{g\in G}$ define
\begin{equation}\label{eq:moment-positive-convolution}
 T_\alpha(p)=\sum_{\beta\le\alpha}\binom\alpha\beta
     c_{\alpha-\beta}(-1)^{|\beta|}
                              \sum_{g\in G}p_g g^\beta.
\end{equation}
The finite programme
\begin{equation}\label{eq:moment-positive-LP}
 \min q\quad\hbox{subject to}\quad
 p_g\ge0,\quad\sum_gp_g=1,\quad q\ge0,\quad
 |T_\alpha(p)-\widetilde y_\alpha|\le q
                                      \quad(|\alpha|\le n)
\end{equation}
has a rational optimizer $p^*$.  If $e_f$ is the variation norm in
\eqref{eq:moment-factor-TV}, its minimum satisfies
\begin{equation}\label{eq:moment-feasible-residual}
 q_*\le e_y+e_f+n\sqrt2h.
\end{equation}
The probability $\widehat\mu_0=\sum_gp_g^*\delta_{Rg}$ obeys
\begin{equation}\label{eq:moment-positive-output-bound}
 W_1(\widehat\mu_0,\mu_0)
 \le R\left\{\frac{C}{m}
            +\Lambda_m(2e_y+e_f+4m\sqrt2h)\right\}.
\end{equation}
Every output atom lies within $C(\rho+Rh)$ of $A_0$ and retains
the canonical coordinates supplied by the geometric stage.
\end{lemma}
\begin{proof}
Round a point $z\in R^{-1}A_0$ to a nearest point of $h\Z^2$,
using a fixed rule in case of a tie.  The rounding distance in
maximum norm is at most $h/2$.  The prior coordinate reserve puts
the rounded point in $[-1/2,1/2]^2$ for small $h$, and
\[
 d_\infty(g,R^{-1}K_A)\le\rho/R+h/2.
\]
It is therefore retained.  Conversely, every retained point has
Euclidean distance at most $C(\rho/R+h)$ from $R^{-1}A_0$.
The rational distance computation in
Lemma~\ref{lem:moment-protected-cutoff} proves that $G$ is a
finite, explicitly decidable set.

Let $U=X/R$, let $\widetilde U$ be the scaled uniform point of
$D$, and let $Z_0=Z/R$.  The exact polygon moments imply
\[
 T_\alpha(p)=\E(\widetilde U-Z_p)^\alpha,
 \qquad \mathcal L(Z_p)=\sum_gp_g\delta_g,
\]
with independent factors.  Thus \eqref{eq:moment-positive-convolution}
has a positive probabilistic realization for every feasible $p$.
For fixed $z$ in the allowed box, the monomial
$(u-z)^\alpha$ is bounded by one.  Replacing $U$ by
$\widetilde U$ changes its expectation, even after averaging in
$z$, by at most $e_f$.  On the same box its gradient with respect
to $z$ has Euclidean norm at most $|\alpha|$.  Quantizing $Z_0$
by the rounding just described therefore changes the convolution
moment by at most $n\sqrt2h$.

The resulting grid law, whose weights need not be rational, is a
comparison point of the real probability simplex.  Its discrepancy
from the rational data is at most $e_y+e_f+n\sqrt2h$, which
proves \eqref{eq:moment-feasible-residual}.  To obtain a rational
optimizer, add the harmless bound
$q\le1+\max_{|\alpha|\le n}|\widetilde y_\alpha|$.
Every $T_\alpha$ has absolute value at most one, so this makes a
nonempty compact rational polytope without changing the minimum.
A linear functional attains its minimum at a vertex, and a vertex
is obtained by solving a nonsingular rational linear system.
The optimum and an optimizing vector are therefore rational and
can be found by finite exact linear programming.  The unknown
comparison law was needed only to bound the optimum.

For the optimizer, the true and fitted convolution moments differ
by at most $q_*+e_y\le2e_y+e_f+n\sqrt2h$.  The two individual
geometric factors have moment discrepancy at most $e_f$.
Lemma~\ref{lem:moment-explicit-conditioning} consequently applies
with $b=2e_y+e_f+n\sqrt2h$.  Scaling transportation distance by
$R$ proves \eqref{eq:moment-positive-output-bound}.  The support
claim follows from the padding estimate already established.
\end{proof}

\begin{corollary}[A sparse positive output]
\label{cor:moment-positive-compression}
The optimizer in Lemma~\ref{lem:moment-positive-programme} can be
chosen with at most
\begin{equation}\label{eq:moment-sparse-output-size}
             \binom{4m+2}{2}=8m^2+6m+1
\end{equation}
nonzero rational weights.  This reduction preserves every fitted
convolution moment through degree $4m$, the objective value, and
the transportation bound.  It uses no further observations.
\end{corollary}
\begin{proof}
There are $d=\binom{4m+2}{2}-1$ nonconstant monomials of total
degree at most $4m$.  Associate to each retained grid point the
rational vector of these $d$ convolution moments together with a
first coordinate equal to one.  If more than $d+1$ weights of a
rational optimizer are positive, the corresponding columns have
a nonzero rational linear dependence $v$.  Its first coordinate
gives $\sum_gv_g=0$, so $v$ has both positive and negative
entries.  Set
\[
       \lambda=\min_{v_g>0}p_g/v_g,
                  \qquad p'_g=p_g-\lambda v_g.
\]
The new weights are nonnegative and rational; their sum and all
the convolution moments are unchanged.  At least one positive
weight vanishes.  Repeating this finite reduction proves the
stated bound.  The hypotheses used in
\eqref{eq:moment-positive-output-bound} are preserved exactly.
\end{proof}

The following proposition retains the all-node observation design.
Its factor, positivity, normalization and error allocations are common
to the default signed-record controller; substituting
Theorem~\ref{thm:linear-moment-sampling} gives
\eqref{eq:two-field-joint-default-cost}, as proved in
Corollary~\ref{cor:linear-joint-budget}.

\begin{proposition}[Complete factorization budget: deterministic alternative]
\label{prop:moment-complete-budget}
The construction of Theorem~\ref{thm:two-field-finite-law} may use
the rational factor, shared acquisition, padded programme, and
sparse output above.  Its stated observation, coordinate, and
transportation bounds are unchanged, and its probability output
can have at most $\binom{4m+2}{2}$ atoms.
Conditional on geometric estimates with the stated bounds, the
entire moment construction remains valid for an arbitrary compact
launch probability.
\end{proposition}
\begin{proof}
Let $a_m$ be the dyadic tolerance in
\eqref{eq:two-field-moment-tolerance}.  Reconstruct geometry at
tolerance $c a_m$ with failure allowance $\delta/2$.  Refine the
rational polygons deterministically so that $\xi+\zeta\le Cc a_m$
and $\rho\le Cc a_m$.  Choose the occupation moment tolerance
$\eta=c a_m$ and the normalized mesh $h$ dyadic and comparable
to $c a_m/m$.  The two factor errors in
\eqref{eq:moment-positive-output-bound} satisfy
\[
 e_y\le C a_m,\qquad e_f\le C a_m,
          \qquad 4m\sqrt2h\le C a_m.
\]
Since $\Lambda_m\le(Cm)^{4m+2}$, taking the fixed constant $C_1$
in \eqref{eq:two-field-moment-tolerance} sufficiently large gives
$C\Lambda_m a_m\le1/m$ for every $m\ge2$.  The positive output
therefore has $W_1$ error at most $C/m$.  With
$m=\lceil C_0/\varepsilon\rceil$, this is the asserted
$C\varepsilon$ error.  The geometry was obtained at the finer
tolerance $Ca_m$.

Apply Lemma~\ref{lem:moment-shared-quadrature} with $n=4m$ and
the remaining failure allowance $\delta/2$.  Its conditional
confidence bound is uniform over every outcome in the geometric
accuracy event.  Thus both stages succeed with probability at
least $1-\delta$, and their observation total is
\[
 N_{\rm geom}(c a_m,\delta/2)
          +Cm^2a_m^{-4}\log\frac{Cm}{a_m\delta}.
\]
Substituting the polynomial geometric bound and
$\log(1/a_m)\le C m\log(Cm)$ gives
\eqref{eq:two-field-joint-exponential-cost}.  All occupation
queries use the original vectors $te_1,-te_1$ and their fixed
prefix translates.  Their dyadic coordinate denominators have
logarithms at most $C m\log(Cm)$; the geometric-stage coordinates
have the same bound.  Charging every attempted bit and command
occurrence therefore gives the stated command-description bound.

The candidate grid has at most $C(m/a_m)^2$ points, whereas
Corollary~\ref{cor:moment-positive-compression} bounds the number
of nonzero output atoms by \eqref{eq:moment-sparse-output-size}.
Their locations and weights are rational and are obtained by
finite exact arithmetic.  Construction of the polygons, moment
integration, linear programming, and weight reduction use no
additional observations.  Their arithmetic work and the
representation lengths of the output weights are separate from
the attempted-bit and command-description bounds.  The moment
proofs used compact support, the component gap, and the geometric
factor estimates; they used no density assumption.  For the joint
uniform experiment, the quantitative density priors enter through
$N_{\rm geom}$ exactly as stated in the geometric theorem.
\end{proof}
